\documentclass[12pt,a4paper]{report}
% Общая преамбула. Подключается из subject/main.tex после \documentclass.

% ============================================================
% Базовые шрифты и кодировка
% ============================================================
\usepackage[T2A]{fontenc}
\usepackage[utf8]{inputenc}
\usepackage[russian]{babel}

\renewcommand{\rmdefault}{cmr}
\renewcommand{\sfdefault}{cmss}
\renewcommand{\ttdefault}{cmtt}
\renewcommand{\familydefault}{\rmdefault}

% ============================================================
% Математика и базовое оформление
% ============================================================
\usepackage{amsmath,amssymb,amsthm,mathtools}
\usepackage{mathrsfs}
\usepackage{bbold}
\usepackage{enumitem}
\usepackage{xparse}
\usepackage{xcolor}
\usepackage{microtype}
\usepackage{geometry}          % для настройки полей
\usepackage{parskip}           % для расстояний между абзацами (без параметров)
\usepackage{hyperref}
\usepackage[nameinlink,noabbrev]{cleveref}
\usepackage{tikz}
\usetikzlibrary{
    positioning,
    arrows.meta,
    calc,
    intersections,
    patterns,
    shapes.geometric,
    decorations.markings,
    matrix,
    fit,
    backgrounds
}

% ============================================================
% Рисунки, графики, схемы, таблицы и фрагменты кода
% ============================================================
\usepackage{graphicx}      % изображения
\usepackage{float}         % точное размещение рисунков/таблиц через [H]
\usepackage{caption}       % подписи к рисункам и таблицам
\usepackage{subcaption}    % несколько рисунков внутри одной фигуры
\usepackage{wrapfig}       % обтекание рисунков текстом
\usepackage{pgfplots}      % графики функций и численных данных
\pgfplotsset{compat=1.18}
\usepackage{tikz-cd}       % коммутативные диаграммы

\usepackage{booktabs}      % аккуратные таблицы
\usepackage{array}
\usepackage{tabularx}
\usepackage{longtable}
\usepackage{multirow}

\usepackage{listings}      % код: C/C++/Python/SQL и др.

% Дополнительная математика
\usepackage{bm}            % жирные математические символы: \bm{x}
\usepackage{bbm}           % красивый blackboard bold для цифр: \mathbbm{0}, \mathbbm{1}
\usepackage{esint}         % \oiint, \varoiint и другие интегралы
\usepackage{cancel}        % \cancel{}, \bcancel{}, \xcancel{}

% ============================================================
% Отступы и расстояния
% ============================================================
\geometry{a4paper,top=2cm,bottom=2cm,left=3cm,right=3cm,marginparwidth=1.75cm}
% parskip уже загружен без параметров, что даёт нулевой отступ первой строки
% и положительный межстрочный интервал между абзацами.

\setlength{\emergencystretch}{3em}

\hypersetup{
    colorlinks=true,
    linkcolor=black,
    urlcolor=black,
    citecolor=black,
    linktoc=all
}

% ============================================================
% Доказательства
% ============================================================
\addto\captionsrussian{%
    \renewcommand{\proofname}{Доказательство}%
}

\makeatletter
\renewenvironment{proof}[1][\proofname]{%
    \par
    \begingroup
    \leftskip=\dimexpr-\@totalleftmargin\relax
    \rightskip=0pt
    \parindent=0pt
    \noindent\textit{\underline{#1.}}\par\nobreak
    \noindent\ignorespaces
}{%
    \unskip\nobreak\hspace{0.15em}\ensuremath{\blacksquare}%
    \par
    \endgroup
}
\makeatother

% ============================================================
% Заголовки
% ============================================================
\usepackage{titlesec}
\titleclass{\chapter}{straight}

\titleformat{\chapter}[display]
  {\normalfont\huge\bfseries\raggedright}
  {\chaptername\ \thechapter}
  {0.7em}
  {\Huge\raggedright}

\titleformat{\section}
  {\normalfont\Large\bfseries}
  {\thesection}{0.7em}{}

\titleformat{\subsection}
  {\normalfont\large\bfseries}
  {\thesubsection}{0.7em}{}

\titleformat{\subsubsection}
  {\normalfont\normalsize\bfseries}
  {\thesubsubsection}{0.7em}{}

\titlespacing*{\chapter}{0pt}{0.45em}{0.45em}
\titlespacing*{\section}{0pt}{0.55em}{0.25em}
\titlespacing*{\subsection}{0pt}{0.45em}{0.18em}
\titlespacing*{\subsubsection}{0pt}{0.35em}{0.15em}

% ============================================================
% Лекции
% ============================================================
\newcounter{lection}
\NewDocumentCommand{\newlection}{m}{%
    \refstepcounter{lection}%
    \par\smallskip
    {\centering\color{gray!70}\normalsize Лекция \thelection. #1\par}
    \vspace{0.1em}
}

% ============================================================
% Теоремы, определения, замечания
% ============================================================
\makeatletter
\NewDocumentCommand{\definition}{o +m}{%
    \par
    \noindent
    \ifdim\@totalleftmargin>0pt
        \hspace*{\dimexpr-\@totalleftmargin\relax}%
    \fi
    \textbf{Определение\IfValueT{#1}{ (#1)}.}~#2
    \par
}
\makeatother

\ExplSyntaxOn

\NewDocumentCommand{\theorem}{o +m}{%
    \par\noindent
    \textbf{%
        Теорема
        \IfValueT{#1}{%
            \regex_match:nnTF { \A \d+ \Z } { #1 }
                {~#1}
                {~(#1)}
        }%
        .
    }~#2
    \par
}

\ExplSyntaxOff

\NewDocumentCommand{\proposal}{o +m}{%
    \par\noindent
    \textbf{Предложение\IfValueT{#1}{~#1}.}~#2
    \par
}

\NewDocumentCommand{\fact}{o +m}{%
    \par\noindent
    \textit{Факт\IfValueT{#1}{ (#1)}.}~#2
    \par
}

\NewDocumentCommand{\question}{+m}{%
    \par\noindent
    \textbf{Вопрос.}~#1
    \par
}

\NewDocumentCommand{\note}{+m}{%
    \par\noindent
    \textbf{Замечание.}~#1
    \par
}

\NewDocumentCommand{\corollary}{o +m}{%
    \par\noindent
    \textbf{Следствие\IfValueT{#1}{~#1}.}~#2
    \par
}

\NewDocumentCommand{\example}{+m}{%
    \par\noindent
    \textbf{Пример.}~#1
    \par
}

% ============================================================
% Списки
% ============================================================
\NewDocumentCommand{\numbers}{+m}{%
    \begin{enumerate}[leftmargin=2em,itemsep=0.25em,topsep=0.3em]
    #1
    \end{enumerate}
}

\NewDocumentCommand{\bullets}{+m}{%
    \begin{itemize}[leftmargin=2em,itemsep=0.25em,topsep=0.3em]
    #1
    \end{itemize}
}

\NewDocumentCommand{\provehere}{+m}{%
    \par\smallskip
    \begin{proof}#1\end{proof}
}

\NewDocumentCommand{\provewthen}{m +m}{%
    \par\smallskip
    \begin{proof}[#1]#2\end{proof}
}

\NewDocumentCommand{\provebullets}{+m}{%
    \begin{proof}
    \begin{itemize}[leftmargin=2em,itemsep=0.25em,topsep=0.3em]
    #1
    \end{itemize}
    \end{proof}
}

\NewDocumentCommand{\provenumbers}{+m}{%
    \begin{proof}
    \begin{enumerate}[leftmargin=2em,itemsep=0.25em,topsep=0.3em]
    #1
    \end{enumerate}
    \end{proof}
}

% ============================================================
% Дополнительные команды
% ============================================================
\usepackage{empheq}      % для \encircle
\usepackage{changepage}  % для \blockindent

% Леммы
\newtheorem{lemma_env}{Лемма}[section]
\newcommand{\lemma}[2][]{\begin{lemma_env}[#1]#2\end{lemma_env}}

% Отступ (переименовано из \indent, чтобы не ломать ядро LaTeX)
\newcommand{\blockindent}[1]{%
  \begin{adjustwidth}{1.5cm}{}#1\end{adjustwidth}%
}

% Одностраничные блоки
\newcommand{\singlepage}[1]{\begin{samepage}#1\end{samepage}}

% Сборки формул
\renewcommand{\gather}[1]{\begin{gather*}#1\end{gather*}}
\renewcommand{\multline}[1]{\begin{multline*}#1\end{multline*}}

% Обрамлённые формулы
\newcommand{\encircle}[1]{\begin{empheq}[box=\fbox]{align*}#1\end{empheq}}

% Математические операторы и символы
\newcommand{\proddots}{\cdot\ldots\cdot}
\newcommand{\abs}[1]{\left|#1\right|}
\newcommand{\angles}[1]{\left\langle#1\right\rangle}
\newcommand{\der}[2]{\frac{\partial #1}{\partial #2}}
\DeclareMathOperator{\differential}{d}
\newcommand{\dif}{\differential\!}
\newcommand{\fmax}{\lor}
\newcommand{\fmin}{\land}
\DeclareMathOperator*{\esssup}{ess\,sup}
\DeclareMathOperator*{\essssup}{ess\,sup}
\DeclareMathOperator*{\esssssup}{ess\,sup}
\newcommand{\Map}{\longrightarrow}
\newcommand{\almost}[1]{\underset{#1}{\overset{\textup{почти всюду}}\Map}}
\newcommand{\circlesign}[1]{\mathrel{\tikz[baseline=-0.5ex]{\node[draw, circle, inner sep=1pt] (char) {\vphantom{$-$}$#1$};}}}

% Дополнительные команды для доказательств
\newcommand{\indentlemma}[3][]{\blockindent{\lemma[#1]{#2\prove[Доказательство леммы]{#3}}}}
\newcommand{\prove}[2][Доказательство]{\begin{proof}[#1]#2\end{proof}}

% ============================================================
% Часто используемые математические обозначения
% ============================================================
\newcommand{\N}{\mathbb{N}}
\newcommand{\Z}{\mathbb{Z}}
\newcommand{\Q}{\mathbb{Q}}
\newcommand{\R}{\mathbb{R}}
\renewcommand{\C}{\mathbb{C}}
\newcommand{\Ord}{\mathrm{Ord}}
\newcommand{\Card}{\mathrm{Card}}

\renewcommand{\o}{\varnothing}
\newcommand{\bs}{\setminus}
\newcommand{\inv}{^{\complement}}
\newcommand{\then}{\Rightarrow}
\newcommand{\map}{\to}
\newcommand{\bydef}{\coloneqq}
\newcommand{\all}[1]{\left\{\begin{aligned}#1\end{aligned}\right.}
\newcommand{\any}[1]{\left[\begin{aligned}#1\end{aligned}\right.}
\newcommand{\exist}{\exists}
\newcommand{\dom}{\operatorname{dom}}
\newcommand{\rng}{\operatorname{rng}}
\newcommand{\defset}[2]{\left\{#1\;\middle|\;#2\right\}}
\newcommand{\switch}[1]{\left\{\begin{aligned}#1\end{aligned}\right.}

% ------------------------------------------------------------
% Большие операторы: индексы автоматически ставятся сверху/снизу
% ------------------------------------------------------------
\let\oldlim\lim
\renewcommand{\lim}{\oldlim\limits}
\let\oldsup\sup
\renewcommand{\sup}{\oldsup\limits}
\let\oldinf\inf
\renewcommand{\inf}{\oldinf\limits}
\let\oldmax\max
\renewcommand{\max}{\oldmax\limits}
\let\oldmin\min
\renewcommand{\min}{\oldmin\limits}

\let\oldsum\sum
\renewcommand{\sum}{\oldsum\limits}
\let\oldprod\prod
\renewcommand{\prod}{\oldprod\limits}
\let\oldcoprod\coprod
\renewcommand{\coprod}{\oldcoprod\limits}

\let\oldbigcap\bigcap
\renewcommand{\bigcap}{\oldbigcap\limits}
\let\oldbigcup\bigcup
\renewcommand{\bigcup}{\oldbigcup\limits}
\let\oldbigsqcup\bigsqcup
\renewcommand{\bigsqcup}{\oldbigsqcup\limits}

\let\oldbigvee\bigvee
\renewcommand{\bigvee}{\oldbigvee\limits}
\let\oldbigwedge\bigwedge
\renewcommand{\bigwedge}{\oldbigwedge\limits}
\let\oldbigoplus\bigoplus
\renewcommand{\bigoplus}{\oldbigoplus\limits}
\let\oldbigotimes\bigotimes
\renewcommand{\bigotimes}{\oldbigotimes\limits}
\let\oldbigodot\bigodot
\renewcommand{\bigodot}{\oldbigodot\limits}
\let\oldbiguplus\biguplus
\renewcommand{\biguplus}{\oldbiguplus\limits}

% Красивые blackboard-bold цифры (\mathbb не во всех шрифтах содержит цифры)
\newcommand{\bbzero}{\mathbbm{0}}
\newcommand{\bbone}{\mathbbm{1}}

% Универсальные обозначения для разных курсов
\newcommand{\norm}[1]{\left\lVert #1 \right\rVert}
\newcommand{\floor}[1]{\left\lfloor #1 \right\rfloor}
\newcommand{\ceil}[1]{\left\lceil #1 \right\rceil}
\newcommand{\paren}[1]{\left( #1 \right)}
\newcommand{\brackets}[1]{\left[ #1 \right]}
\newcommand{\set}[1]{\left\{ #1 \right\}}
\newcommand{\conj}[1]{\overline{#1}}
\newcommand{\vect}[1]{\bm{#1}}

\DeclareMathOperator{\RePart}{Re}
\DeclareMathOperator{\ImPart}{Im}
\DeclareMathOperator{\Arg}{Arg}
\DeclareMathOperator{\Ker}{Ker}
\DeclareMathOperator{\Image}{Im}
\DeclareMathOperator{\rank}{rank}
\DeclareMathOperator{\tr}{tr}
\DeclareMathOperator{\sgn}{sgn}
\DeclareMathOperator{\lcm}{lcm}
\DeclareMathOperator{\Span}{span}

% Differential notation in math; preserve the text dot-below accent.
\let\textdotbelow\d
\DeclareRobustCommand{\d}{\ifmmode\mathrm{d}\else\expandafter\textdotbelow\fi}

\endinput

% Всё ниже оставлено из исходного файла как справочный образец титульной
% страницы и LaTeX не читает после команды \endinput.
% ============================================================
% Информация о документе и титульная страница
% ============================================================
\date{III семестр, 2026 г.}
\title{Название предмета. Неофициальный конспект}
\author{Лектор: ФИО лектора\\
Конспект: Викторова Дарья, 25.Б03-пу}

\begin{document}

\begin{titlepage}
    \thispagestyle{empty}
    \begin{center}
        \vspace*{0.6cm}

        {\large Санкт-Петербургский государственный университет\par}
        \vspace{0.35em}
        {\normalsize Факультет прикладной математики~--- процессов управления\par}

        \vspace*{5.1cm}

        {\fontsize{30}{35}\selectfont\bfseries Название предмета\par}
        \vspace{0.65em}
        {\fontsize{17}{21}\selectfont\itshape Неофициальный конспект\par}

        \vspace{2.7cm}

        {\normalsize
        \begin{tabular}{@{}rl@{}}
            \textit{Лектор:} & ФИО лектора \\
            \textit{Конспект:} & Викторова Дарья, 25.Б03-пу
        \end{tabular}\par}

        \vfill

        {\large III семестр\par}
        \vspace{0.22em}
        {\large 2026 год\par}
        \vspace*{0.55cm}
    \end{center}
\end{titlepage}

\tableofcontents
\newpage

\setcounter{lection}{0}
\newlection{Число Месяц 2026 г.}

\end{document}



\title{Теоретическая механика. Неофициальный конспект}
\author{Лектор: ФИО лектора\\
Конспект: Викторова Дарья, 25.Б03-пу}
\date{III семестр, 2026 г.}

\begin{document}

\begin{titlepage}
    \thispagestyle{empty}
    \begin{center}
        \vspace*{0.6cm}

        {\large Санкт-Петербургский государственный университет\par}
        \vspace{0.35em}
        {\normalsize Факультет прикладной математики~--- процессов управления\par}

        \vspace*{5.1cm}

        {\fontsize{30}{35}\selectfont\bfseries Теоретическая механика\par}
        \vspace{0.65em}
        {\fontsize{17}{21}\selectfont\itshape Неофициальный конспект\par}

        \vspace{2.7cm}

        {\normalsize
        \begin{tabular}{@{}rl@{}}
            \textit{Лектор:} & Потоцкая Ирина Юрьевна \\
            \textit{Конспект:} & Викторова Дарья, 25.Б03-пу
        \end{tabular}\par}

        \vfill

        {\large III семестр\par}
        \vspace{0.22em}
        {\large 2026 год\par}
        \vspace*{0.55cm}
    \end{center}
\end{titlepage}

\tableofcontents
\newpage

\setcounter{lection}{0}

\newlection{1 сентября 2026 г.}

\section*{Литература}
\begin{enumerate}
    \item \textbf{Бабаджанянц Л.\ К. <<Классическая механика>>. --- СПб. : Соло, 2007}.
    \item Бухгольц Н.\ Н. <<Основной курс теоретической механики>>. В 2 т. --- М. : Наука, 1966.
    \item Маркеев А.\ П. <<Теоретическая механика>>. --- М. : Наука, 1990.
    \item Поляков Н.\ Н. <<Теоретическая механика>>. --- М. : Высшая школа, 1975.
    \item Мещерский И.\ В. <<Сборник задач по теоретической механике>>. --- М. : Наука, 1970.
    \item \textbf{Бать М.\ И., Джанелидзе Г.\ Ю., Кельзон А.\ С. <<Теоретическая механика в примерах и задачах>>. В 3 т. --- М. : Наука, 1972}.
\end{enumerate}

\medskip

\textit{Цель теоретической механики} – найти закон изменения состояния объекта (зависимость состояния объекта от времени).

% ============================================================
% Схема 1: закреплённая (не плавающая)
% ============================================================
\begin{center}
\begin{tikzpicture}[
    emptybox/.style={draw=none, rectangle, thick, align=center, text width=6.0cm, inner sep=4pt, font=\normalsize},
    mybox/.style={draw, rectangle, thick, align=center, minimum width=2.5cm, inner sep=4pt, font=\normalsize},
    midbox/.style={draw, rectangle, thick, align=center, text width=6.0cm, inner sep=6pt, font=\normalsize},
    bigbox/.style={draw, rectangle, thick, align=left, text width=6.0cm, inner sep=6pt, font=\normalsize},
    arrow/.style={->, thick},
    node distance=1.2cm
]

\node (goal) [emptybox] {Цель};
\node (tasks) [emptybox, below=1.2cm of goal] {\underline{Задачи}};

% Левая ветвь
\node (study) [midbox, below=1.2cm of tasks, xshift=-4.5cm] 
    {Изучение процесса движения,\\ состояния объекта};

\node (chars) [bigbox, below=1.2cm of study] 
    {Нахождение характеристик движения. \\
    --- Кинематика (геометрическое движение) \\
    $\bullet$ Скорость \\
    $\bullet$ Ускорение \\
    $\bullet$ Траектория \\
    --- Динамика \\
    $\bullet$ Количество движения (импульс движения) \\
    $\bullet$ Кинетическая энергия \\
    $\bullet$ Кинетический момент (момент количества движения) \\
    --- Статика \\
    $\bullet$ Равновесие систем и сил
};

% Правая ветвь
\node (model) [mybox, below=1.2cm of tasks, xshift=4.5cm] 
    {Составление математической модели};

\node (solver) [bigbox, below=1.2cm of model] 
    {1. Выбрать систему координат (задать координаты) \\
    2. Составить уравнение движения относительно заданной системы координат \\
    3. Решить уравнение движения
};

% Стрелки
\draw[arrow] (goal) -- (tasks);
\draw[arrow] (tasks.south) -- (study.north);
\draw[arrow] (study) -- (chars);
\draw[arrow] (tasks.south) -- (model.north);
\draw[arrow] (model) -- (solver);

\end{tikzpicture}
\end{center}

% ============================================================
% Примеры математических моделей
% ============================================================
\subsection{Примеры математических моделей}

1) Робот-пылесос

% ============================================================
% Схема 2: закреплённая (не плавающая)
% ============================================================
\begin{center}
\begin{tikzpicture}[
    scale=1.2,
    every node/.style={font=\normalsize},
    >=stealth,
    thick
]
% Параметры
\def\radius{2.2}
\def\wheelwidth{0.8}
\def\wheelheight{1.2}

\coordinate (center) at (0,0);

% Левое колесо
\draw[thick] (-\radius-\wheelwidth, -\wheelheight/2) rectangle (-\radius, \wheelheight/2);
% Правое колесо
\draw[thick] (\radius, -\wheelheight/2) rectangle (\radius+\wheelwidth, \wheelheight/2);

% Окружность
\draw[thick] (center) circle(\radius);

% Диаметр
\draw[thin] (-\radius,0) -- (\radius,0);
% l
\node[below] at (0,-0.3) {$l$};

% Векторы скоростей колёс
\node[above] at (-\radius-\wheelwidth/2, \wheelheight/2+0.3) {$\vec V_L$};
\node[above] at (\radius+\wheelwidth/2, \wheelheight/2+0.3) {$\vec V_R$};

% Вектор V из центра вверх
\draw[->, thick] (center) -- (0, \radius+1.0) node[above] {$\vec V$};

% Вектор Omega
\draw[->, thick] (0.6,0.4) arc[start angle=-20, end angle=20, radius=0.6] node[midway, right] {$\vec \Omega$};

% Стрелка на правом колесе и подпись r
\draw[<->, thick] (\radius+\wheelwidth/2, 0) -- (\radius+\wheelwidth/2, \wheelheight/2) node[midway, right] {$r$};

\end{tikzpicture}
\end{center}

\medskip
\noindent
Введём обозначения для скоростей колёс:
\[
\vec V_R = r\,\vec\omega_R, \qquad \vec V_L = r\,\vec\omega_L .
\]

Тогда:
\begin{equation}
\left\{
\begin{aligned}
\vec V &= \frac{\vec V_R + \vec V_L}{2} 
       = \frac{r(\vec\omega_R + \vec\omega_L)}{2} 
       &&\text{— линейное движение},\\[1.2ex]
\vec\Omega &= \frac{r(\vec\omega_R - \vec\omega_L)}{l},\quad \omega_R > \omega_L 
       &&\text{— поворот вокруг центра масс}.
\end{aligned}
\right.
\label{eq:robot}
\end{equation}

\medskip
Рассмотрим два частных случая:

\begin{enumerate}
\item[а)] $\vec\Omega = 0 \;\Rightarrow$ движение прямо $\;\Rightarrow\; 
   \vec\omega_R = \vec\omega_L = \vec\omega \neq 0$.  
   Тогда линейная скорость:
   \[
   \vec V = \frac{r\cdot 2\vec\omega}{2} = r\vec\omega .
   \]

\item[б)] $\vec V = 0 \;\Rightarrow\; \vec\omega_R = -\vec\omega_L$.  
   Тогда угловая скорость поворота:
   \[
   \vec\Omega = \frac{r(\vec\omega_R + \vec\omega_R)}{l} 
              = \frac{2r\,\omega_R}{l}.
   \]
\end{enumerate}

% ============================================================
% 2) Манипулятор
% ============================================================
2) Манипулятор

\end{document}


